% =====================================================================
% Recuperação Híbrida e Raciocínio Multi-hop
% Responsável: Pessoa 4
%
% O que cobrir:
% - Busca vetorial (índices chunk_embedding e entidade_embedding, prefixo 'query: ').
% - Text-to-Cypher: como o LLM gera a consulta, como ela é validada antes de executar.
% - Como a busca vetorial e a consulta Cypher são combinadas (ordem, pesos, fusão de resultados).
% - Travessia multi-hop: exemplo de pergunta, a consulta Cypher gerada e o caminho percorrido.
% - Objeto de resposta enviado à interface (resposta, nós, arestas, Cypher, confiança).
%
% Dicas: use \rot{Pessoa} para rótulos e \rel{CRIOU} para relações.
% Cite com \cite{chave} (chaves em referencias.bib). Ao terminar,
% apague o \pendente abaixo.
% =====================================================================

\section{Recuperação Híbrida e Raciocínio Multi-hop}
\label{sec:recuperacao}

\pendente{recuperação híbrida e raciocínio multi-hop}
